\cotefile{156-4}
\scanpage{50}
\mspage{49}

\deleted{\uline{\unsure{Cont|Compt|Const} 5} Si $F$ et $G$ sont deux figures disjointes, et \unsure{si} $F'$, \deleted{\unsure{$G'$}\illegible{1} de} subdivision\unsure{s|{}}\unsure{,} \inserted{de $F$} alors $F'$ et $G$ sont disjointes.} \margin{\deleted{\uline{NB} $F'\cap G$ est défini \unsure{comme|com.}, \unsure{c'est} $F'$ et $G$ \uline{compatibles} qui est \unsure{en cause}.}}

\deleted{\uline{Cor} Si $F'$ \unsure{raffine|subfine} $F$, $G'$ \unsure{raffine} $G$, \deleted{\unsure{alors}} $F$ et $G$ disj. alors $F'$ et $G'$ disjointes.}

\uline{C\deleted{\unsure{ont|ompt|onst}} 5} Soient $F$ une figure, $F'$ une subdivision de $F$. \deleted{\unsure{A}}Alors pour \unsure{tt} figure $G$, $G$ est disjointe \unsure{;|:|à} $F$ ssi elle est disjoint\unsure{e} de $F'$

\uline{NB} Il suffit de le poser pour $G$ élémentaire.

Il résulte alors de \deleted{\unsure{Cont|Compt|Const} 4} C4 que $F'\cup G$ est une subdivision de $F\cup G$.

\uline{Cor} Soient $F$ et $G$ figures disjointes, alors pour $F'\ll F$, $G'\ll F$, $F'$ et $G'$ sont disjointes, et $F'\cup G'\ll F\cup G$. Si $F'$ et $G'$ sont des \unsure{subd} de $F,G$, alors $F'\cup G'$ en est une de $F\cup G$.

\uline{C 6} (Subdivisions induites) Soient $F$ une figure, $F'$ une subdivision de $F$, $G$ une sous-figure de $F$. Alors il existe une unique subdivision $G'$ de $G$, telle que $G'\subset F'$ (subdivision induite \inserted{et $G'$ est formée des $X\in F'$ telles que $X\ll G$})

\margin{NB \inserted{Double} Transitivité de l'induction. \uline{Notation} $F'|G$}

\uline{C 7} (\unsure{Recollement|Récollement} des subdivisions) Soit $F$ une figure, \deleted{\unsure{chaque} \illegible{3} \unsure{tt}} \inserted{$F_i$ des sous-figures} \deleted{$X\in$} telles que $F=\bigcup F_i$. Soit, pour \unsure{tt} $i$, $F'_i$ une subdivision de $F_i$, et supposons

\scanpage{51}
\mspage{50}

que \unsure{pour} $i,j\in I$, $\unsure{i\Rightarrow j|i\ne j}$, $F'_i|F_i\cap\deleted{G}F_j=\deleted{G}F'_j|F_i\cap F_j$ Alors il existe une unique subdivision $F'$ de $F$ telle que $F'|F_i=F'_i\quad\forall i\in I$

\uline{Cor} Pour que deux subdivisions \inserted{$F',F''$ de $F$} soient égales, il \unsure{f} et \unsure{s} qu'elles induisent la $=$ subdivision sur chaque sous-figure \unsure{élémentaire} $F_{\unsure{X|x}}$ \unsure{.ie|i.e.} sur chaque \unsure{multistrate}.

\deleted{\unsure{Cor}} \uline{Proposition} \deleted{\unsure{1|a}}) Soient $F$ une figure, $G$ une sous-figure, $\deleted{G}F'$ un raffinement de $F$, et $G'=\{X\in F'\mid X\ll G\}$ c'est donc une sous figure de $F'$, \deleted{et} \inserted{Je dis qu'}elle raffine $G$ et c'est \inserted{\deleted{\unsure{d'ailleurs|même}}} la plus grande sous-figure de $F'$ qui raffine $G'$.

% slanted
\margin{\uline{raffinement} \uline{induit}. \uline{NB} Si $\deleted{G}F'\subset\deleted{G}F$, alors $G'$ n'est autre que l'intersection $F'\cap G'$.}

\diagram{\begin{tikzcd}[ampersand replacement=\&,column sep=small,row sep=small] \deleted{G}F' \arrow[r,phantom,"\ll" description] \arrow[d,phantom,"\cup" description] \& F \arrow[d,phantom,"\cup" description] \\ G' \arrow[r,phantom,"\ll" description] \& G \end{tikzcd}}

\deleted{\unsure{Ceci dit},} Il y a qque chose à prouver, car le seul fait que les $X\in G'$ raffinent $G$ n'implique \unsure{ps|pas} que $G'$ raffine $G$. Mais on a une \deleted{\unsure{raffinement}} \inserted{subdiv.} $\overline{F\unsure{''|'|{}}}$ de $F$ tel que \diagram{\begin{tikzcd}[ampersand replacement=\&,column sep=small,row sep=small] F' \arrow[r,phantom,"\subset" description] \& \overline{F''}\arrow[r,"\sim","s."' ]\arrow[d,phantom,"\cup" description] \& F\arrow[d,phantom,"\cup" description] \\ \& \overline{G''}\arrow[r,"\sim","s"' ] \& G \end{tikzcd}} \unsure{considérons|considérer} la subd. induite \deleted{sur} $G$\inserted{(\unsure{C,6|C6})}, $\overline{G''}$, formée des $X\in\overline{F}$ \unsure{telles|tels} que $X\ll \deleted{\unsure{F}}G$. \deleted{\unsure{Ceci}} Par définition, on a donc $G'=F'\cap G'$, qui est \deleted{bien} une sous-figure de $\overline{\deleted{\unsure{C}}G''}$ donc raffine $G$.

\scanpage{52}
\mspage{51}

\uline{Corollaire} Sous les conditions précédentes, $G'$ est le $\mathrm{Inf}$ de $F'$ et $G$, \unsure{dans|ds} l'\unsure{ens} des figures pour la \deleted{\unsure{relation|notion}} \inserted{relation} de raffinement. En d'autres termes, \deleted{\unsure{toute|tt}} figure \inserted{$\unsure{\Phi|Q|\varphi}$} qui raffine à la fois $F'$ et $\deleted{\unsure{F}}G$, raffine $G'$

\diagram{\begin{tikzcd}[ampersand replacement=\&,column sep=small,row sep=small,execute at end picture={\draw (\tikzcdmatrixname-1-1.west) -- (\tikzcdmatrixname-1-2.north east); \draw (\tikzcdmatrixname-2-3.south west) -- (\tikzcdmatrixname-1-4.north east); \draw (\tikzcdmatrixname-2-4.south west) -- (\tikzcdmatrixname-1-5.north east); \draw (\tikzcdmatrixname-2-5.south west) -- (\tikzcdmatrixname-1-6.west);}] \deleted{\Phi} \& \deleted{\unsure{\Phi\subset\Phi|G\subset\Phi}} \& \deleted{\overline{F}'} \arrow[r,"\deleted{\sim}"] \arrow[d,phantom,"\deleted{\cup}" description] \& \deleted{F'} \arrow[r,phantom,"\deleted{\subset}" description] \arrow[d,phantom,"\deleted{\cup}" description] \& \deleted{\overline F} \arrow[r,"\deleted{\sim}"] \arrow[d,phantom,"\deleted{\cup}" description] \& \deleted{F} \arrow[d,phantom,"\deleted{\cup}" description] \\ \& \& \deleted{\overline G'} \arrow[r,"\deleted{\sim}"] \& \deleted{G'} \arrow[r,phantom,"\deleted{\subset}" description] \& \deleted{\overline G} \arrow[r,"\deleted{\sim}"] \& \deleted{G} \end{tikzcd}}

On utilise que $\unsure{F|F'}\subset\overline{F''}$, où $F''$ subd. de $F$, et que $\Phi\subset\overline{F} '$, où $\overline F'$ subd. de $F'$, \deleted{et que}

\diagram{\begin{tikzcd}[ampersand replacement=\&,column sep=small,row sep=small] \& \& \deleted{\widetilde F}\arrow[d,"\deleted{\unsure{\subset|\swarrow}}" description,phantom]\arrow[dr,"\sim"] \& \& \\ \Phi\arrow[r,phantom,"\subset" description] \& \overline F'\arrow[r,dash,"\sim"]\arrow[d,phantom,"\cup" description] \& F'\arrow[r,phantom,"\subset" description]\arrow[d,phantom,"\cup" description] \& \unsure{\overrightarrow F|\overline F}\arrow[r,"\sim"]\arrow[d,phantom,"\cup" description] \& F\arrow[d,phantom,"\cup" description] \\ \& \overline G'\arrow[r,"\sim"] \& G'\arrow[r,phantom,"\subset" description] \& \overline G\arrow[r,"\sim"] \& G \end{tikzcd}}

On va prouver que $\Phi\subset\overline G'$\inserted{(\deleted{\unsure{donc|dont}} $\Phi\ll G'$)}. Soit $X\in\Phi$, on a $X\in\overline F'$, \deleted{et $X$} il faut prouver $X\in\overline{G'}$, \unsure{pour|i.e.|or} \deleted{$\unsure{\overline{F}'|G'}$} que \unsure{tt} élément de $\overline{F'}$ \uline{qui raffine} $G$, est \unsure{dans|ds} $\overline G'$. Mais $\overline F'$ est \uline{induite par une} subd. $\widetilde F$ de $F$ (C 4 revu et corrigé -- cf son cor.) donc $\overline F'\subset\widetilde F\xrightarrow{\sim}\overrightarrow F$, donc $X\in\widetilde F$ et $X$ raffine $G$, donc \deleted{$X\in\unsure{\widetilde F'|\overline F'}$} $X\in\widetilde G$, subd. de $G$ induite par la \unsure{subd.} $\widetilde F$ de $F$, donc

\diagram{\begin{tikzcd}[ampersand replacement=\&,column sep=small,row sep=small] \overline F'\arrow[r,phantom,"\subset" description] \& \widetilde F\arrow[r,"\sim"]\arrow[d,phantom,"\cup" description] \& \overrightarrow F\arrow[r,"\sim"]\arrow[d,phantom,"\cup" description] \& F\arrow[d,phantom,"\cup" description] \\ \& \widetilde G\arrow[r,"\sim"] \& \overline G\arrow[r,"\sim"] \& G \end{tikzcd}}

\deleted{\unsure{$\Phi$}} $X\in\overline F'\cap\widetilde G=\overline G'$, qed, ouf ! (C'était l'ancien \unsure{Cont|Const|Compt} 8)

\diagram{\begin{tikzcd}[ampersand replacement=\&,column sep=small,row sep=small,execute at end picture={\draw ([xshift=-3pt]\tikzcdmatrixname-4-1.west) .. controls ([xshift=-16pt,yshift=43pt]\tikzcdmatrixname-2-1.west) and ([yshift=14pt]\tikzcdmatrixname-1-2.north) .. ([yshift=10pt]\tikzcdmatrixname-1-3.north) .. controls ([xshift=38pt,yshift=10pt]\tikzcdmatrixname-1-5.north) and ([xshift=15pt]\tikzcdmatrixname-4-5.east) .. ([xshift=4pt]\tikzcdmatrixname-5-5.south east);}] \deleted{\unsure{F'|G'}}\arrow[r,dash] \& \deleted{F'} \& \widetilde F\arrow[dr,"\sim"]\arrow[d,dash] \& \& \vphantom{F} \\ \Phi\arrow[r,hook] \& \overline F'\arrow[ur,hook]\arrow[dr,"\sim"]\arrow[dd,hook] \& {} \& \overrightarrow F\arrow[dr,"\sim"]\arrow[dd,hook] \& \\ \& \& F'\arrow[ur,hook]\arrow[d,hook]\arrow[dd,dash,bend right=13] \& \& F\arrow[d,hook] \\ \vphantom{G} \& \overline G'\arrow[dr,"\sim"]\arrow[r,phantom,"\subset" description] \& \widetilde G\arrow[r,"\sim"] \& \overline G\arrow[r,"\sim"] \& G \\ \& \& G'\arrow[ur,phantom,"\subset" description] \& \& \vphantom{G} \end{tikzcd}}

\scanpage{53}
\mspage{52}

\uline{13 juin} On n'a pas encore exprimé tous les Axiomes \uline{\unsure{Cont|Contr}}\unsure{.} 1, \unsure{....|...}, il \unsure{reste} \unsure{pour compte|par exemple}

\uline{Cont} 3 \unsure{dim} finie

\uline{Cont} 4 ``admissibilité'' -- mais résulte déjà je crois des axiomes C.

\uline{Cont} 6 \inserted{Critère de} disjonction \unsure{via les lieux} -- on y reviendra.

\uline{Cont} 7 Deux lieux de \unsure{str.} dominantes différentes sont disjoints.

Mais avant de continuer, je voudrais encore une fois changer d'optique !

Prendre comme \unsure{ens} de base $\unsure{\mathcal F|\mathscr F}$, l'ens des figures, avec les deux relations d'ordre \[ \leq\ \text{(ou l'inclusion)} \qquad \text{et}\qquad \unsure{\preceq|\preccurlyeq}\ \text{(la subdivision)} \]

\uline{Contrée} (quatrième mouture !)

\unsure{$U$ est|C'est} un ens. $\mathcal F$ d'objets, appelés \uline{figures}, et deux relations d'ordre $\leq$ et $\preceq$ sur $\mathcal F$ : \[ \unsure{\mathcal C|\mathscr C|\mathcal E}=(\unsure{\underline{\mathcal F}|\mathcal F},\leq,\preceq) \] satisfaisant à des axiomes\unsure{,|{}} qu'on va énumérer. Si $F$ et $G$ sont deux figures et $F\leq G$, on dit que $F$ est une \uuline{sous}-\uline{figure}, et on \unsure{écrira|écrit} aussi souvent $F\subset G$. Si $F\preceq G$, on dit que $F$ est une \uline{subdivision} de $G$, ou que $F$ subdivise $G$.

\scanpage{54}
\mspage{53}

D'abord, on va dégager les axiomes pour $\leq$

\uline{$\unsure{\mathcal C|\mathscr C|\mathcal E}$ 1} a) $\unsure{\mathcal F|\mathscr F}\ne\varnothing$

\deleted{\uline{\unsure{C 2}}} b) Si $F\in\mathcal F_{\leq}$, et $(F_i)_{i\in I}$ est une famille d'objets de $\mathcal F$ avec $F_i\leq F$, alors $\underset{i}{\mathrm{Sup}}\,F_{\unsure{j|i}}$ existe. (Toute famille majorée de figures a une borne sup.)

\uline{Corollaire} $\mathcal F$ a un plus petit élé\deleted{\unsure{ment|mt}} (savoir, le $\mathrm{Sup}$ de la famille vide d'objets)

\uline{NB} En présence de \deleted{\unsure{C2}} \uline{C 1 b)}, la validité de a) \uline{équivaut} à celle du corollaire.

Le plus petit \unsure{élémt|élément} est appelé \uline{figure vide}, \unsure{notée|notation} $\varnothing_{\mathcal F}$ ou simplement $\varnothing$, par abus de langage, et si confusion n'est à craindre

\deleted{\unsure{Réun.|Fam.|Réun} \unsure{des|de} figures \unsure{comp.}}

\deleted{\uline{Définition} \uline{figure irréductible} $F$ : telle que si $F=\mathrm{Sup}$}

Figures \uline{compatibles}, ens de figures \deleted{\unsure{compatibles}} compatibles \deleted{\unsure{: un qui}} \inserted{dans leur ens.} (ou \uline{configuration}) : qui est majoré pour $\leq$. \margin{\unsure{?|P}} Figure engendrée, ou \uline{réunion} ($F\vee G$, $\bigvee_i\unsure{f|F}_i$, $\mathrm{Sup}\,F_i$, parfois noté $F\cup G$ ou $\deleted{\unsure{F}}\bigcup_iF_i$ cf plus bas).

\uline{Déf} \uline{Figure irréductible} $F$ : telle que $\displaystyle F=\underset{i\in I}{\mathrm{Sup}}\,F_i\ \overset{\unsure{-}}{\Longrightarrow}\ \exists i\in I,\ F_i=F$.

\deleted{\unsure{C 2}} Les figures irréductibles sont appelées des \uline{multistrates}. \inserted{L'ens des multistrates noté $\unsure{\mathcal M|\mathscr M}$} Les figures irréductibles

\scanpage{55}
\mspage{54}

majorées par $F$ s'appellent les \uline{strates} de $\deleted{\unsure{E}}F$. (Donc, \deleted{\unsure{de|des}} strate de $F$ est une multistrate)

L'ens des strates de $F$ \deleted{s'appelle} sera noté $\overline F$ : \[ \begin{aligned} \deleted{\unsure{\overline F\subset\mathcal M|F\in\mathcal M}} &\qquad \unsure{\mathcal M|\mathscr M}\subset\unsure{\mathcal F|\mathscr F}\\ F\longrightarrow\overline F &\qquad \mathcal F\longrightarrow\unsure{\mathcal P|\mathscr P}(\mathcal M). \end{aligned} \]

\uline{$\unsure{\mathcal C|\mathscr C|\mathcal E}$ 2} Toute figure est la borne supérieure de ses strates : \[ F=\underset{\deleted{\unsure{A}}X\in\overline F}{\mathrm{Sup}}\ X \]

\uline{Corollaire} 1 L'application $\deleted{X}F\longrightarrow\overline{\deleted{\unsure{C|G}}F}$ \[ \mathcal F\longrightarrow\mathcal P(\mathcal M) \] est \unsure{injection|injective}, et induit un iso d'\unsure{ens} ordonnés de $\mathcal F$ avec l'ens $\unsure{\Phi|Q|\varphi}\subset\mathcal P(\mathcal M)$ formé des $\overline F$.

Cet iso. transforme $\underset{i}{\mathrm{Sup}}$ en $\bigcup_i$

\uline{Corollaire 2} L'ens $\Phi$ des parties de $\mathcal M$ \unsure{satisfait|satisfaisant|satisfont} les deux conditions

\deleted{a) $\Phi\ne\varnothing$}

% slanted
\margin{\deleted{\unsure{Ceci|Ceci ?}}} \deleted{\unsure{b)}}a) Si $A\in\Phi$, et $(A_i)_{i\in\unsure{I|\Phi}}$ est une famille \unsure{di'|d'}éléments de $\Phi$, avec $A_i\subset A$, alors $\bigcup_iA_i\in\Phi$.

b) $\forall X\in\mathcal M$, il existe un plus petit $A=A_X$ tel que $X\in\unsure{A|A_x|A_X}$ (a fortiori, $\Phi$ recouvre $\mathcal M$).

\scanpage{56}
\mspage{55}

c) $\varnothing\in\unsure{\Phi|Q|\varphi}$ \deleted{\unsure{$\ne\varnothing$} [en vertu de b), c'est tjrs vrai si $\unsure{\mathcal M|\mathscr M}\ne\varnothing$], et \unsure{si} $\mathcal M=\varnothing$, \unsure{cela signifie que} \unsure{$\varnothing$}}

\deleted{\uline{\unsure{Rép.|Rq.}} (équivaut : $\varnothing\in\Phi$)}

\uline{Remarque} Inversement, soient $\mathcal M$ un ens, $\deleted{\Phi}\unsure{\mathcal F|\mathscr F}\subset\unsure{\mathcal P|\mathscr P}(\mathcal M)$ satisfaisant les conditions a) b) c). NB La condition c résulte de \inserted{a),} b si $\mathcal M\ne\varnothing$. Dans tous les cas, elle équivaut à : \deleted{\unsure{$\in$|à}} c') $\deleted{\Phi}\mathcal F\ne\varnothing$ \unsure{i.e} existe $\deleted{X}A\subset\mathcal M$ telle que $A\in\deleted{\Phi}\mathcal F$

% slanted
\margin{\unsure{Ces|Les} conditions impliquent que $\mathcal F$ est stable par $A\cap B$ (mais en général on n'a \unsure{$\not\Rightarrow$|pas} $\mathcal M\in\mathcal F$ !)}

Alors

1°) $\Phi$ ordonné par inclusion \deleted{\unsure{ad|sa}}satisfait les conditions C1, C2. \deleted{Les objets irréductibles} \deleted{\unsure{sont}}

2°) Les objets irréductibles sont les $A_X$, et $X\mapsto\deleted{\unsure{X}}A_X$ est une bijection de $\mathcal M$ \deleted{avec} l'ens des objets \unsure{irréd.} de $\deleted{\Phi}\mathcal F$.

3°) Si $X\in\mathcal M$, $A\in\deleted{\Phi}\mathcal F$, alors $X\in A\Longleftrightarrow A_X\subset A$ \deleted{, donc} (tautologique) donc \deleted{\unsure{$A\in\mathcal F$}} si on identifie $\mathcal M$ \inserted{(par $X\unsure{\vdash|\mapsto} A_X$)} à l'ens des objets irréductibles de $\deleted{\Phi}\mathcal F$, alors pour $A\in\mathcal F$, $\overline A=\{A_X\mid A_X\subset A\inserted{\text{\unsure{ie.}}\,X\in A}\}$ s'identifie à $A$ lui-$=$.

\centerline{\rule{1.8cm}{0.35pt}}

On \unsure{considérera|considère} $\mathcal M$ comme ordonné par $\leq$ (jamais par $\unsure{\prec|\preceq}$ !) : multistrates et \uline{sous-multistrates}. Les figures \inserted{$F$} peuvent s'identifier aux \deleted{\unsure{des parties}} ens. de multistrates $\overline F$ correspondants (chaque figure à l'ens de ses strates).

\scanpage{57}
\mspage{56}

\uline{Proposition} a) Si $F$ \deleted{\unsure{est|était}} est une figure, alors les figures $F'$ telles que $F'\leq F$ s'identifient, par l'application $F'\mapsto\overline F'$, aux parties de $\overline F$ qui sont \uline{fermées} pour la relation d'ordre \inserted{$\leq$} de $\unsure{\mathcal M|\mathscr M}$.

b) Si $F$ et $G$ sont deux figures, $F\wedge G$ existe (inf pour $\leq$) \inserted{\uuline{figure intersection}} et \deleted{\unsure{on|en}} $\overline{F\cap G}=\overline F\cap\overline G$. Même pour une famille $(F_i)_{i\in I}$ \unsure{quelc.|qelc.} \inserted{avec \uuline{$I\neq\varnothing$}}.

% margin note written sideways
\margin{\uline{Mais attention} il vaut mieux ne pas les noter $F\cap G$, car quand il y a une notion de support, ça ne correspond pas à l'intersection des supports, sauf si $F,G$ compatibles.}

[ \uline{Lieux} de $\unsure{\mathcal F|\mathscr F}$ : ce sont les \deleted{\unsure{pas|pet}} figures qui sont \uline{minimales} dans $\mathcal F\setminus\{\varnothing_{\mathcal F}\}$. \deleted{\unsure{Un|Une},} \deleted{Un lieux est de \unsure{ce nom|nom}}  Tout lieu est irréductible. Donc les lieux sont aussi les éléments minimaux de $\mathcal M$ : ce sont les multistrates qui n'ont d'autres strates \unsure{qu'elles-même|qu'elles-mêmes|qu'elle-même}. ] \deleted{Si $F$ est une figure, les}

% slanted
\margin{Définition \unsure{courante|courte} n'est raisonnable que dans certains cas \deleted{\unsure{(quand les|(quand la}} (p. ex. les contrées \unsure{en boules''|en boules}, ou contrées \unsure{par morceaux|ponctuées} $\ldots$) \unsure{Car dans|Car chez} les \unsure{cas|cat.} types d'\unsure{ordre|données} \unsure{quelconq.|quelconques}, \unsure{un|une} \unsure{modèle|modélé} $\overline E$ \unsure{de str.|de structure} \unsure{semi-}\illegible{1} \inserted{\unsure{univers vides|univers vide},} forme \deleted{\unsure{peut être une|peut être un}} \inserted{une} multistrate \unsure{minimale|minime}, mais on ne \unsure{l'appelle|appelle} un ``lieu'' que si elle est \unsure{réduite|discrète|vide réelle} !}

\uline{Figures} \uline{disjointes} \unsure{,,|''} \uline{$F$} et \uline{$G$} : telles que a) $F$ et $G$ compatibles, i.e $F\cup G$ existe \deleted{et b)} ou encore $\{F,G\}$ majoré pour $\leq$ et b) $F\cap G=\varnothing_{\mathcal F}$ \deleted{\unsure{ou|et} Cela signifie \unsure{donc|que}} \inserted{(i.e. $\overline F\cap\overline G=\varnothing$)}.

% slanted
\margin{\unsure{bloque|bloquer} mes ``\unsure{coproduits|composantes}''}

\deleted{Si $F$ et $G$ sont} \uline{Prop} Si $F$ et $G$ sont disjointes, alors pour $F'\leq F$, $G'\leq G$, $\unsure{F'|F'_{a}}$ et $G'$ sont disjointes, et $F'\cup G'$ est une sous figure de $F\cup G$, l'application $(F';G')\mapsto F'\cup G'$ est une bijection \deleted{de l'\unsure{ens|ens.} de $\unsure{\mathcal F|F}$}
